% ID: SEL-D0039-OA-E04
% ESTADO: cotejada. Segunda pasada visual y borrador compilado revisados.
% Fuente: Selectividad/matematicas_aplicadas/Ua4MATEMÁTICAS APLICADAS examen.pdf
% SHA-256: 78ff6bdd48cbac71e6cc47d6ddf3dcb9492540ecb9cbe2032ebfec02e39fb66d
% Página: 1; opción A; ejercicio original 4.

La calificación en Matemáticas de los alumnos de un centro docente es una variable aleatoria que sigue una distribución Normal de desviación típica 1.2. Una muestra de 10 alumnos ha dado las siguientes calificaciones:
\[
3\quad8\quad6\quad3\quad9\quad1\quad7\quad7\quad5\quad6.
\]
\begin{enumerate}[label=\alph*)]
\item \textbf{(1.75 puntos)} Se tiene la creencia de que la calificación media de los alumnos del centro en Matemáticas es a lo sumo 5 puntos. Con un nivel de significación del 5\%, plantee el contraste unilateral correspondiente $(H_0:\mu\leq5)$, determine la región crítica y razone si la creencia es fundada o no.
\item \textbf{(0.75 puntos)} \textquestiondown{}Obtendría la misma respuesta si el nivel de significación fuese del 15\%?
\end{enumerate}
